\chapter{CMake 依赖、安装边界和 CI}

\section{问题入口：项目在别人机器上构建不了}

本机能构建不代表项目健康。常见问题包括：头文件路径靠 IDE 隐式配置，库路径写死在本机目录，测试命令只存在个人脚本里，依赖版本没有记录。工程化的目标是让项目在新机器上用少量命令复现。

\section{目标级依赖}

CMake 里应围绕 target 建模：

\begin{lstlisting}[language=CMake,caption={目标级依赖}]
add_library(config_loader src/config_loader.cpp)
target_compile_features(config_loader PUBLIC cxx_std_20)
target_include_directories(config_loader PUBLIC include)

add_executable(config_loader_tests tests/config_loader_tests.cpp)
target_link_libraries(config_loader_tests PRIVATE config_loader)
\end{lstlisting}

\code{PUBLIC} 表示使用者也需要这个属性，比如公开头文件路径。\code{PRIVATE} 表示只对当前目标有效。不要全局 \code{include_directories}，它会让模块意外依赖不该看到的路径。

\section{库目标和可执行目标分开}

如果所有逻辑都写在 \code{main.cpp}，测试很难复用。把核心逻辑做成库：

\begin{lstlisting}[language=CMake,caption={库和程序分开}]
add_library(wordfreq_core
    src/split_words.cpp
    src/count_words.cpp
    src/top_words.cpp)
target_include_directories(wordfreq_core PUBLIC include)
target_compile_features(wordfreq_core PUBLIC cxx_std_20)

add_executable(wordfreq src/main.cpp)
target_link_libraries(wordfreq PRIVATE wordfreq_core)

add_executable(wordfreq_tests tests/wordfreq_tests.cpp)
target_link_libraries(wordfreq_tests PRIVATE wordfreq_core)
\end{lstlisting}

程序和测试都链接同一个核心库。这样测试覆盖的是生产代码，而不是复制出来的一份逻辑。

\section{依赖版本要明确}

如果项目需要第三方库，可以用系统包、子模块、包管理器或 CMake FetchContent。无论方式如何，都要记录版本。不要让“最新版本”成为隐式依赖。今天能构建，半年后最新版本变了，行为也可能变。

对小项目，尽量先用标准库。引入依赖要回答：

\begin{itemize}
  \item 它解决的问题是否足够大？
  \item 是否有稳定版本和许可证？
  \item 构建和交叉平台成本是多少？
  \item 如果依赖停止维护，替换成本多高？
\end{itemize}

\section{安装和导出不是第一天必需，但边界要清楚}

库如果要给其他项目使用，需要安装头文件和库文件：

\begin{lstlisting}[language=CMake,caption={安装目标示意}]
install(TARGETS wordfreq_core
        ARCHIVE DESTINATION lib
        LIBRARY DESTINATION lib
        RUNTIME DESTINATION bin)

install(DIRECTORY include/ DESTINATION include)
\end{lstlisting}

教学项目不一定要完整安装规则，但目录结构应该为未来留余地：公开头文件在 \code{include}，内部实现不暴露。

\section{CI 最小流程}

持续集成不是大公司专属。最小 CI 只需要在每次提交后运行：

\begin{enumerate}
  \item 配置项目。
  \item 编译。
  \item 运行测试。
  \item 可选：运行格式检查、静态分析、消毒器构建。
\end{enumerate}

伪命令：

\begin{lstlisting}[language=bash,caption={CI 核心命令}]
cmake -S . -B build -DCMAKE_BUILD_TYPE=Debug
cmake --build build --parallel
ctest --test-dir build --output-on-failure
\end{lstlisting}

如果这些命令不能在干净环境运行，说明项目还有隐式依赖。

\section{构建类型矩阵}

常见矩阵：

\begin{itemize}
  \item Debug：快速编译，适合单元测试。
  \item Release：发现优化下的问题，跑性能 smoke test。
  \item ASan/UBSan：发现内存和未定义行为。
  \item TSan：发现数据竞争，适合并发模块。
\end{itemize}

不一定每次提交都跑全部矩阵。可以把快速测试放在每次提交，慢测试放在夜间或发布前。关键是规则明确。

\section{构建失败也要可诊断}

CI 失败时，输出要足够定位问题。测试名字不要叫 \code{test1}；断言失败要能说明输入和期望；构建脚本不要吞掉错误。一个好的失败信息应该让维护者知道下一步看哪个模块。

\begin{exercisebox}
给配置加载器项目写一个 CMake 结构：核心库、命令行程序、测试目标。要求测试只链接核心库，不运行命令行程序。再设计一个 CI 命令清单，包含 Debug 测试和 ASan 测试。
\end{exercisebox}
